\documentclass[11pt,a4paper]{article}
\usepackage[T1]{fontenc}
\usepackage[utf8]{inputenc}
\usepackage{lmodern,microtype}
\usepackage[margin=.85in]{geometry}
\usepackage{amsmath,amssymb,amsthm,mathtools,aliascnt}
\usepackage{booktabs,tabularx,array,enumitem,needspace,float}
\usepackage[section]{placeins}
\usepackage{xcolor,tikz,mdframed}
\usetikzlibrary{arrows.meta,positioning,calc,fit}
\usepackage[hidelinks]{hyperref}
\usepackage[nameinlink,noabbrev]{cleveref}
\usepackage{fancyhdr}
\hypersetup{pdftitle={Sync HotStuff RR v2},pdfsubject={L1-scheduled leaders and generation-wide recovery}}
\setlength{\headheight}{14pt}
\pagestyle{fancy}\fancyhf{}
\lhead{\small Sync HotStuff RR v2}\rhead{\small L1-scheduled leaders}
\cfoot{\thepage}
\setlength{\parindent}{0pt}\setlength{\parskip}{4pt}
\setlength{\emergencystretch}{2em}
\setlist{nosep,leftmargin=*}
\newcolumntype{Y}{>{\raggedright\arraybackslash}X}
\definecolor{rrRed}{HTML}{A51C30}
\definecolor{rrBlue}{HTML}{2563EB}
\definecolor{rrTeal}{HTML}{0F766E}
\definecolor{rrAmber}{HTML}{B45309}
\newcommand{\RR}[1]{{\color{rrRed}#1}}
\newcommand{\QC}{\mathcal C}
\newcommand{\DC}{\mathcal F}
\newcommand{\CB}{\mathsf{CB}}
\newcommand{\Hash}{\mathsf H}
\newcommand{\rank}{\operatorname{rank}}
\newcommand{\tagg}[1]{\mathsf{#1}}
\newtheorem{definition}{Definition}[section]
\newaliascnt{assumption}{definition}
\newtheorem{assumption}[assumption]{Assumption}
\aliascntresetthe{assumption}
\newaliascnt{lemma}{definition}
\newtheorem{lemma}[lemma]{Lemma}
\aliascntresetthe{lemma}
\newaliascnt{theorem}{definition}
\newtheorem{theorem}[theorem]{Theorem}
\aliascntresetthe{theorem}
\newaliascnt{proposition}{definition}
\newtheorem{proposition}[proposition]{Proposition}
\aliascntresetthe{proposition}
\crefname{assumption}{assumption}{assumptions}
\Crefname{assumption}{Assumption}{Assumptions}
\crefname{lemma}{lemma}{lemmas}
\Crefname{lemma}{Lemma}{Lemmas}
\crefname{theorem}{theorem}{theorems}
\Crefname{theorem}{Theorem}{Theorems}
\newcounter{protocol}
\newenvironment{protocol}[2]{\refstepcounter{protocol}\label{#1}
\begin{mdframed}[nobreak=true,linewidth=.6pt,skipabove=8pt,skipbelow=8pt,
innertopmargin=7pt,innerbottommargin=7pt]\small
\textbf{Protocol \theprotocol. #2}\par\smallskip}{\end{mdframed}}
\crefname{protocol}{Protocol}{Protocols}
\Crefname{protocol}{Protocol}{Protocols}
\title{Sync HotStuff RR v2\\\large L1-Scheduled Leaders and Generation-Wide Recovery}
\author{}\date{11 October 2026}
\begin{document}
\maketitle
\begin{abstract}
We extend standard-synchronous Sync HotStuff with a finalized L1 leader
schedule, ordinary leader terms, and a generation-wide failure boundary.
One valid quit proof freezes all unfinished timers in its generation;
every correct replica recovers to the same view $(g+1,0)$ without Byzantine
agreement on the triggering proof. Recovery retains Sync HotStuff's entry
locks, two waiting stages, leader-selected NEWVIEW, and anchor re-vote.
Regular successor votes require a terminal commit certificate. A successor
may distribute its proposal earlier, but cannot obtain admitted votes before
the predecessor's terminal commitment. This gate prevents repeated faulty
successors from starving finality. We give execution diagrams, safety and
liveness proofs, and communication and latency bounds. Red protocol text
marks extensions to the original standard-synchronous protocol.
\end{abstract}
The schedule is available before execution; L1 provides common leader
identities, not a per-recovery consensus service. The construction uses
ordinary signatures and $f+1$-vote certificates, as in Section II of
Sync HotStuff~\cite{sync}. It is a commit-gated reference version.
Immediate terminal-QC rotation is discussed separately in
\cref{sec:optimization}.

\section{State machine replication}
\label{sec:smr}
A deterministic state machine executes a sequence of client commands.
Replicas order commands in hash-linked blocks; committing a block commits
its uncommitted ancestors in order. We use the following SMR contract.
\begin{definition}[SMR]
An execution satisfies \emph{safety} if the committed logs of any two correct
replicas are prefix-comparable and no completed commitment is withdrawn.
It satisfies \emph{validity} if committed blocks contain only admissible
commands and valid state transitions. It satisfies \emph{request liveness}
if every well-formed correct-client request that remains admissible for log
inclusion until ordered is eventually committed by every correct replica.
\end{definition}
Committed-height progress is a separate intermediate property: every
correct replica eventually commits blocks of arbitrarily large height.
Request liveness additionally requires the fair input-service assumption
below. A vote certificate alone is not a finality certificate.

\section{Model and security assumptions}
\label{sec:model}
\begin{assumption}[Authenticated synchronous committee]
\label{ass:network}
The fixed committee has $n=2f+1$ replicas. At most $f$ are Byzantine;
the remaining replicas retain their state and remain correct throughout
the execution. Complete messages between correct replicas, including all
data and evidence needed for validation, are delivered and processed within
a known $\Delta>0$. Clocks run at the same rate; offsets need not agree.
Signatures are unforgeable and $\Hash$ is collision resistant.
There are finitely many events in every bounded time interval.
At an exact timer boundary, process received evidence before callbacks.
\end{assumption}
The complete-message premise includes dependencies that are not already
cached. It is an execution assumption, not a claim that arbitrarily large
histories can be carried by a capacity-limited network at constant cost.
The communication analysis charges the evidence actually transmitted.

\begin{assumption}[Finalized L1 schedule]
\label{ass:l1}
L1 authenticates one immutable committee configuration and a deterministic
leader function $L(g,t)$ for $g,t\in\mathbb N$. The applicable schedule and
its verification material are available to every correct replica before
the corresponding view can be entered. Every scheduled leader belongs to
the fixed committee. There is no L1 transaction on the block or recovery
critical path. L1 does not choose a recovery anchor or certify an SMR commit.
\end{assumption}
For height liveness, the sequence of initial leaders must revisit correct replicas:
for every $g_0$ there is a finite $g\ge g_0$ with correct $L(g,0)$.
A cyclic enumeration of the committee across generations is one concrete
schedule satisfying this condition. For request liveness, also require that
for every fixed $g$ and every $t_0$ there is a finite $t\ge t_0$ with correct
$L(g,t)$. Cyclic term orders satisfy this second condition. Safety requires
neither fairness property and does not require a scheduled leader to be correct.

\begin{assumption}[Initialization and fair input service]
\label{ass:input}
Correct replicas share a committed genesis block $B_0$ and enter the initial
view $(0,0)$ within $\Delta$ of one another. Correct leaders can produce
valid empty blocks. A correct client's request is eventually disseminated
to all correct replicas. Requests have a common admission order in which
each request has finitely many predecessors. Every correct leader fills its
fresh blocks with the oldest pending admissible requests that fit, and a
request is removed from the pending queue only after commitment. Requests
already present in the proposed ancestry need not be duplicated: committing
the child also commits those ancestors. Each admitted request fits in one
batch. Correct replicas retain and serve committed data and evidence.
\end{assumption}

\section{Views, certificates, and retained evidence}
\label{sec:notation}
A view is $v=(g,t)$, ordered lexicographically. Generation $g$ changes only
on failure recovery; term $t$ changes on regular rotation. Its unique leader
is $L(v)=L(g,t)$. Every signed message binds the protocol/configuration
domain, message type, generation, term, height, block hash, opening identity,
and parent where applicable. Signatures from one generation cannot be
reused in another. A correct replica signs at most one DATA block at each
$(g,t,k)$; redelivery cannot generate another vote or restart its timer.

A fresh block has the form $B_k=(x_k,\Hash(B_{k-1}))$ and must satisfy the
application validity predicate. A DATA vote certificate
$\QC_v(B_k)$ contains $f+1$ matching votes by distinct members. An ANCHOR
vote is explicitly tagged and may certify the unchanged recovery anchor;
its certificate is usable as the first fresh child's parent justification.
Both types are ranked by $\rank(\QC_v(B_k))=(g,t,k)$.
Different blocks with the same rank have a fixed hash tie-break when a
highest certificate must be selected. Correct replicas never decrease the
rank of a retained entry lock across recoveries. The lock for the current
generation's ANCHOR first vote is fixed when STATUS is sent; later arrivals
can update the retained highest QC but do not retroactively raise that
generation's entry lock.

After recovery, the scheduled initial leader announces an opening block
$A_g$ in a signed NEWVIEW. Its identity is $(g,h(A_g),\Hash(A_g))$;
different encodings or different vote subsets for the same block do not
create different openings. Term $t$ owns the $K\ge1$ fresh heights
\[
 h(A_g)+tK+1,\ldots,h(A_g)+(t+1)K.
\]
The opening is initially certified, not assumed committed. Its ANCHOR
vote does not consume one of the $K$ fresh positions and starts no commit
timer. Genesis is the common initial opening for generation zero.

\subsection{Terminal commitment and normal admission}
\label{sec:done}
Let $T_{g,t}$ be the block at height $h(A_g)+(t+1)K$ and
$N_{g,t}=\QC_{g,t}(T_{g,t})$ its terminal DATA QC. A correct replica
multicasts
\[
 \langle\tagg{DONE},g,t,h(T_{g,t}),\Hash(T_{g,t}),
             \Hash(A_g)\rangle_r
\]
only after the original $2\Delta$ timer for $T_{g,t}$ successfully commits
it. A terminal commit certificate $\DC_{g,t}(T_{g,t})$ comprises $f+1$
distinct matching DONE signatures, the terminal QC, and the complete
admission history. Receiving an imported commit certificate does not
authorize a new DONE signature. A correct recipient may commit the
certificate's block and ancestors because at least one signer is correct
and actually completed that terminal timer.

Normal entry to $(g,t+1)$ requires $\DC_{g,t}(T_{g,t})$, not just $N_{g,t}$.
Entry to a later normal term requires every intervening terminal commit
certificate. Each DATA proposal binds its opening identity; its bundle
contains the signed opening, its ANCHOR QC (except at genesis), the normal
entry certificates, and the parent justification. Cached identical
dependencies may be referenced, provided they are available within the
complete-message bound.

The successor may pre-send its first signed proposal once $N_{g,t}$ is
available. Before $\DC_{g,t}$ arrives, this is speculative data: it gives no
permission to enter the term, vote, start a timer, or quit because of that
successor's proposed fork. Its already-admitted ancestor history is still
checked immediately. On obtaining $\DC_{g,t}$, recipients can validate the
cached proposal and any prospective fork at the now-admissible slot.
Thus the gate constrains voting, not the earliest transmission of bytes.

\subsection{Opening evidence and quit proofs}
\label{sec:evidence}
For $g>0$, an entry witness $E_g$ is any valid quit proof for generation
$g-1$, with its admission history. All such witnesses name the same target
$(g,0)$ regardless of their old term. The signed opening has the form
\[
 R_g=\langle\tagg{NEWVIEW},(g,0),A_g,Q_g,E_g\rangle_{L(g,0)},
\]
where $Q_g$ certifies $A_g$ in a strictly older generation. A common
synthetic genesis certificate has rank $(-1,0,0)$ and can be used when no
non-genesis QC exists. Opening generation zero requires only the common
initial configuration; its base and genesis parent justification are fixed.
Different witnesses or QC encodings for the same opening block do not
define different openings.

Admissibility is checked against the authenticated schedule and complete
entry history, independently of a recipient's branch preference. Term zero
is admissible on its entry witness even if its leader is silent. A DATA
history additionally needs an opening ANCHOR QC. Terms $t>0$ need the
intervening terminal commit certificates. A complete quit proof
$\CB(g,t)$ is one of:
\begin{enumerate}[label=(\alph*)]
\item $f+1$ distinct signatures on $\langle\tagg{BLAME},g,t\rangle$
for an admissible term;
\item two different admissible DATA proposals at the same $(g,t,k)$,
signed by $L(g,t)$, with their complete ancestry;
\item for $g>0$, two valid leader-signed NEWVIEW messages selecting
different opening blocks for $(g,0)$.
\end{enumerate}
Case (c) also covers comparable but different opening blocks: they give
different quota origins. With one opening, an incompatible DATA history has
a first divergent height owned by one scheduled leader, yielding case (b).
Signing the same block twice or changing its certificate encoding is not
equivocation. A single BLAME signature is not a complete quit proof.

A replica freezes a generation at most once. It retains all evidence and
completed commitments. Old-generation proofs are archival and cannot
advance the current generation. A proof for a not-yet-entered generation
is buffered, then processed immediately on authenticated entry. A complete
admissible higher-term proof can freeze its generation even before the
recipient has locally reached that term; its admission certificates already
prove every required preceding terminal commitment.

\section{Protocols}
\label{sec:protocol}
We keep the two-figure organization of the original paper. Black text
retains its propose/vote/timer and lock-based recovery operations. Red text
specifies the composite view, ordinary rotation, public commit gate,
generation-wide evidence handling, and fixed recovery target.

\begin{protocol}{alg:steady}{Steady state and ordinary leader rotation}
Let $v=(g,t)$ be the local active view and $L=L(g,t)$.
\RR{Run DATA operations only while generation $g$ is unfrozen.}
\begin{enumerate}[label=\textbf{\arabic*)},itemsep=5pt]
\item \textbf{Propose.} If $r=L$, upon obtaining
$\QC_v(B_{k-1})$, broadcast
$\langle\tagg{PROPOSE},B_k,v,\QC_v(B_{k-1})\rangle_L$,
where $B_k$ extends $B_{k-1}$.
\RR{Bind the opening identity and its admission bundle. Use the previous
term's terminal QC for an ordinary term's first fresh child; generation
zero starts from genesis and a recovered term zero starts from its ANCHOR
QC. Propose only the $K$ assigned fresh positions. A scheduled successor
may pre-send its first child on the terminal QC, but it remains speculative
until the predecessor's terminal commit certificate is available.}
\item \textbf{Vote.} On a valid proposal, if no leader equivocation is known,
forward the proposal and multicast
$\langle\tagg{DATA\text{-}VOTE},B_k,v\rangle_r$; start
$\tagg{commit\text{-}timer}_{v,k}=2\Delta$ on that first vote.
\RR{First validate the complete admission and compare retained signed
openings and histories, including older terms of $g$. Relay a complete
fork proof and freeze before another vote. Catch up through all intervening
terminal commit certificates before voting. Do not DATA-vote before the
predecessor DONE certificate admits the term. If already beyond $B_k$
on the same certified history, relay its QC without re-voting. A certified
child with a valid opening ANCHOR QC is processable even if the recipient
declined that opening's original ANCHOR vote; the personal entry-lock
comparison applies to ANCHOR voting, not to this child.}
\item \textbf{(Non-blocking) Commit.} When the timer reaches zero, commit
$B_k$ and its uncommitted ancestors.
\RR{Process evidence at that instant first. Commit only if the timer's
generation has not frozen; normal term changes leave old timers running.
On successful expiration of the terminal block's own timer, multicast
DONE once. Assemble and relay $\DC_{g,t}$, whose committed prefix may be
installed on receipt. If unfrozen, enter $(g,t+1)$ on this certificate,
clear the pending local complaint, and allow the buffered successor
proposal to trigger a DATA vote. Disable the completed term's progress
test; retain its equivocation listener. Imported commits never generate
DONE signatures, and completed commits are never withdrawn.}
\end{enumerate}
\end{protocol}

\begin{protocol}{alg:recovery}{Generation exit and SH-style recovery}
Let $v=(g,t)$ be a current-generation view.
\begin{enumerate}[label=\textbf{\roman*)},itemsep=5pt]
\item \textbf{Blame and quit.} If fewer than $p$ fresh proposals trigger
DATA votes after $(2p+4)\Delta$ from local entry, broadcast
$\langle\tagg{BLAME},v\rangle_r$.
On gathering $f+1$ matching signatures, broadcast the certificate.
On leader equivocation, broadcast the original signed witnesses.
\RR{Schedule a separate check for every $p\in\{1,\ldots,K\}$ at the
local deadline $\operatorname{entry}(v)+(2p+4)\Delta$. Keep at most one
pending local complaint. Count distinct admitted fresh positions for which
the replica has either issued a DATA vote or obtained a valid DATA QC
through authenticated catch-up. Unvoted speculative proposals and ANCHOR
votes do not count. A complete admissible current-generation
quit proof, even for an older or locally unreached term, freezes all of $g$.
Stop new DATA and ANCHOR votes and normal entries; cancel every unfinished
generation-$g$ commit timer. Keep receiving/relaying old votes and forming
QCs. Record certified committed prefixes without reopening the generation.
Do not restart recovery on additional same-generation proofs.}
\item \textbf{Status and entry.} Wait $\Delta$ from local exit. Pick a
highest valid known QC, set the entry lock to it, and send STATUS to the
new leader.
\RR{Every current-generation quit targets $(g+1,0)$, whose unique leader
is the pre-finalized $L(g+1,0)$, regardless of the blamed term. At
$e_r=a_r+\Delta$, set the lock to the higher-ranked of the retained lock
and all known QCs from entered generations through $g$; send the QC and
full verifying history, with a valid quit proof for $g$ serving as
$E_{g+1}$, to that leader.
Enter $(g+1,0)$, clear the local complaint, and reset its proposal clock.
Future-generation messages remain buffered until their authenticated entry.}
\item \textbf{New-view.} The new leader waits $2\Delta$ from its own entry
and broadcasts a signed NEWVIEW containing a highest valid known QC.
\RR{Choose only a QC from a strictly older generation, or genesis.
Pin its block as the one opening $A_{g+1}$ and attach the entry witness.
Two different opening blocks signed for $(g+1,0)$ are equivocation,
including comparable blocks. No BA or per-recovery L1 transaction is run.}
\item \textbf{First vote.} Upon a valid NEWVIEW whose QC rank is at least
the entry-lock rank, forward it and vote for its unchanged anchor.
\RR{Tag this as ANCHOR-VOTE in $(g+1,0)$, at most once, and start no commit
timer. After its ANCHOR QC forms, the leader proposes the first fresh child
and resumes \cref{alg:steady}. Before any vote, inspect retained openings
and history for equivocation. Late old-generation evidence cannot freeze
$g+1$. During entry, first process a buffered complete current-generation
quit proof, if present; this may immediately freeze the new generation.}
\end{enumerate}
\end{protocol}

\subsection{Processing rules that matter to the proofs}
\label{sec:processing}
A valid QC and a locally acceptable ANCHOR first vote are different notions.
The entry lock may prevent a replica from voting for an opening. If a
valid ANCHOR QC nonetheless forms, a later certified child has a new-view
parent QC and must be processed, unless signed equivocation is known.
Requiring every DATA voter to have personally accepted the original
NEWVIEW would break the timer proof. Conversely, a correct initial leader
chooses a QC at least as high as every correct entry lock, so all correct
replicas vote for its opening.

Before using a history for catch-up, preserve and compare its signed
opening and admitted prefixes. Two roots are compared before ordinary
slot allocation; with the same root, compare the first divergent DATA
position. Invalid or incomplete future-term admission is buffered or
rejected, not upgraded to a quit proof. The already-admitted portion of
such a message may still expose genuine old-term equivocation.
Correct replicas relay a newly accepted proposal on their first DATA vote,
and relay an available QC when their certified progress has passed that
proposal. Commit certificates embedded in any admissible history are
recorded even if an accompanying quit proof immediately freezes the node.

For communication bounds, deduplicate by semantic identity, not by byte
encoding. Forward each valid NEWVIEW opening $(g,h(A),\Hash(A))$ once,
regardless of its witness or QC representation. DATA proposals and QCs
are keyed by their message/certificate type, view, height, and block hash;
different ranked QCs for the same block remain distinct. DONE is keyed
by its view, terminal block, and opening. Retain at most two different
signed proposals per admitted slot and two different openings per
generation as fork witnesses. Relay the first complete quit proof, the
first usable QC for each such key, and the first certificate of each
terminal commit; later equivalent representations are stored without
causing unbounded relay traffic. This rule never suppresses the
timer-specific QC required for catch-up or commitment protection.
The bounds count admitted protocol traffic, not arbitrary Byzantine
flooding; fixed-size valid batches and finite evidence-processing capacity
are required by \cref{ass:network}.

\section{Execution diagrams}
\label{sec:figures}
The diagrams separate DATA certification, timed commitment, and ANCHOR
re-certification. Arrows are representative transmissions. Only the
explicitly labeled timing example uses equal actual one-hop delay $\delta$;
the proofs use the worst-case bound $\Delta$.
\begin{figure}[htbp]\centering
% Native TikZ fragment. Illustrative equal-delay execution, not a timing premise.
% Width 15.8 cm; all visible text uses at least 9 pt.
\begin{tikzpicture}[
  x=1cm,y=-1cm,font=\fontsize{9.5}{11}\selectfont,
  lane/.style={draw=black!25,densely dashed},
  msg/.style={-{Stealth[length=1.8mm]},semithick,draw=rrBlue},
  lab/.style={fill=white,inner sep=2pt,align=center},
  state/.style={draw=rrBlue!60,fill=white,rounded corners=2pt,
    align=center,inner sep=4pt,font=\fontsize{9.5}{11}\selectfont},
  timed/.style={draw=rrTeal,dashed,rounded corners=2pt,
    fill=rrTeal!4,align=center,inner sep=4pt},
  phase/.style={anchor=west,font=\fontsize{9.5}{11}\selectfont\bfseries}
]
  \path[use as bounding box] (0,-.5) rectangle (15.8,10.1);
  \node[align=center] at (2.0,0) {Correct leader\\$L(g,t)$};
  \node[align=center] at (7.5,0) {Correct replica\\$P$};
  \node[align=center] at (13.4,0) {Correct replica\\$Q$};
  \foreach \x in {2,7.5,13.4} \draw[lane] (\x,.5)--(\x,9.4);

  \node[phase] at (.1,.85) {Fresh blocks pipeline on QCs};
  \node[lab,anchor=east] at (1.7,1.45) {$0$};
  \draw[msg] (2,1.45)--node[lab] {$\mathrm{PROPOSE}(X)$} (7.5,2.05);
  \draw[msg] (2,1.45)--(13.4,2.05);
  \node[lab,anchor=west] at (13.7,2.05) {$\delta$};
  \node[state,minimum width=8.3cm] at (10.1,2.7)
    {Validate and forward $X$; broadcast DATA votes.\\
     Start each local $2\Delta$ commit timer.};
  \draw[msg] (7.5,3.3)--node[lab] {DATA votes for $X$} (2,3.9);
  \draw[msg] (13.4,3.3)--(2,3.9);
  \node[state] at (2.0,4.5)
    {$2\delta$: form $\QC_{g,t}(X)$;\\propose child $Y$};
  \draw[msg] (2,5.1)--node[lab] {$\mathrm{PROPOSE}(Y,\QC_{g,t}(X))$} (7.5,5.7);
  \draw[msg] (2,5.1)--(13.4,5.7);
  \node[state,minimum width=8.3cm] at (10.1,6.35)
    {$3\delta$: validate and vote for $Y$;\\start another $2\Delta$ timer.};
  \node[timed,minimum width=12.7cm] at (7.9,7.35)
    {The timer for $X$ keeps running while $Y$ and later blocks are certified.};
  \node[state,draw=rrTeal,minimum width=12.7cm] at (7.9,8.4)
    {$9\delta=\delta+2\Delta$: $P$ and $Q$ commit $X$ if generation $g$ is not frozen.\\
     $11\delta=3\delta+2\Delta$: they similarly commit $Y$ and its ancestors.};
  \node[anchor=west,align=left,font=\fontsize{9}{10.5}\selectfont] at (.1,9.6)
    {Illustration: $\delta=\Delta/4$, with equal actual delays. Time runs downward, not to scale.\\
     Arrows represent all-to-all voting; a QC has $f+1$ distinct signers.};
\end{tikzpicture}

\caption{Certification and commitment within one admitted term. A parent QC enables a new child while older DATA timers continue. The example uses actual delay $\delta=\Delta/4$.}
\label{fig:steady}
\end{figure}

\begin{figure}[htbp]\centering
% Native TikZ fragment. The public DONE certificate gates votes, not proposals.
\begin{tikzpicture}[
  x=1cm,y=-1cm,font=\fontsize{9.5}{11}\selectfont,
  lane/.style={draw=black!25,densely dashed},
  msg/.style={-{Stealth[length=1.8mm]},semithick,draw=rrBlue},
  lab/.style={fill=white,inner sep=2pt,align=center},
  state/.style={draw=rrBlue!65,fill=white,rounded corners=2pt,
    align=center,inner sep=4pt},
  gate/.style={draw=rrTeal,fill=rrTeal!5,rounded corners=2pt,
    align=center,inner sep=5pt},
  phase/.style={anchor=west,font=\fontsize{9.5}{11}\selectfont\bfseries}
]
  \path[use as bounding box] (0,-.5) rectangle (15.8,11.4);
  \node[align=center] at (2.1,0) {Current leader\\$L(g,t)$};
  \node[align=center] at (7.8,0) {Correct replica\\$P$};
  \node[align=center] at (13.5,0) {Successor leader\\$L(g,t+1)$};
  \foreach \x in {2.1,7.8,13.5} \draw[lane] (\x,.5)--(\x,10.6);

  \node[phase] at (.1,.85) {Terminal certification may trigger an early proposal};
  \draw[msg] (2.1,1.5)--node[lab] {terminal proposal $T$} (7.8,2.0);
  \draw[msg] (2.1,1.5)--(13.5,2.0);
  \node[state,minimum width=8.9cm] at (10.2,2.65)
    {Vote for $T$ in $(g,t)$ and start its $2\Delta$ timer.\\
     Terminal QC $N_{g,t}=\QC_{g,t}(T)$ forms.};
  \draw[msg] (7.8,3.25)--node[lab] {$N_{g,t}$} (13.5,3.65);
  \draw[msg] (13.5,4.1)--node[lab] {optional early $\mathrm{PROPOSE}(Y,N_{g,t})$} (7.8,4.55);
  \node[state,draw=rrAmber,minimum width=11.8cm] at (7.9,5.25)
    {Buffer $Y$. A terminal QC alone does \textbf{not} permit a DATA vote in term $t+1$.\\
     The old terminal commit timer continues; generation $g$ is unchanged.};

  \node[phase] at (.1,6.1) {Terminal commitment supplies the voting gate};
  \node[gate,minimum width=11.8cm] at (7.9,6.8)
    {After the original terminal timer completes, broadcast $\mathrm{DONE}(g,t,T)$.\\
     Collect $f+1$ matching DONE signatures into $\DC_{g,t}(T)$.};
  \draw[msg,draw=rrTeal] (7.8,7.55)--node[lab] {$\DC_{g,t}(T)$ + complete history} (13.5,8.0);
  \node[gate,minimum width=11.8cm] at (7.9,8.7)
    {Verify $\DC_{g,t}(T)$ and import its committed prefix. Enter $(g,t+1)$.\\
     Now validate and vote for buffered $Y$; start $Y$'s $2\Delta$ timer.};
  \draw[msg] (7.8,9.45)--node[lab] {DATA votes for $Y$; normal pipeline resumes} (13.5,9.9);
  \node[anchor=west,align=left,font=\fontsize{9}{10.5}\selectfont] at (.1,10.85)
    {No L1 transaction and no recovery on this path. Each replica applies the same gate locally.\\
     DONE reports an original timed terminal commit; QC receipt or imported commit does not issue DONE.};
\end{tikzpicture}

\caption{Commit-gated regular rotation. A successor may pre-send its first proposal, but DATA votes wait for the predecessor terminal DONE certificate. This path does not change generation.}
\label{fig:regular}
\end{figure}

\begin{figure}[htbp]\centering
% Native TikZ fragment. Representative messages; time is not drawn to scale.
\begin{tikzpicture}[
  x=1cm,y=-1cm,font=\fontsize{9.5}{11}\selectfont,
  lane/.style={draw=black!25,densely dashed},
  msg/.style={-{Stealth[length=1.8mm]},semithick,draw=rrBlue},
  lab/.style={fill=white,inner sep=2pt,align=center},
  state/.style={draw=rrBlue!65,fill=white,rounded corners=2pt,
    align=center,inner sep=4pt,font=\fontsize{9.5}{11}\selectfont},
  wait/.style={draw=rrAmber,dashed,fill=rrAmber!4,rounded corners=2pt,
    align=center,inner sep=4pt},
  phase/.style={anchor=west,font=\fontsize{9.5}{11}\selectfont\bfseries}
]
  \path[use as bounding box] (0,-.5) rectangle (15.8,17.4);
  \node[align=center] at (1.5,0) {Old leader\\$L(g,t)$};
  \node[align=center] at (5.7,0) {Correct replica $P$\\active at $(g,t_P)$};
  \node[align=center] at (9.9,0) {Correct replica $Q$\\active at $(g,t_Q)$};
  \node[align=center] at (14.1,0) {New leader\\$L'=L(g+1,0)$};
  \foreach \x in {1.5,5.7,9.9,14.1} \draw[lane] (\x,.6)--(\x,16.55);

  \node[phase] at (.1,.95) {(i) Quit the generation, regardless of the active term};
  \node[state,draw=rrRed,minimum width=8.0cm] at (5.0,1.7)
    {At $a_P$, $P$ receives complete $\CB(g,t^*)$.\\
     Relay proof; freeze $g$ and cancel its pending timers.};
  \draw[msg,draw=rrRed] (5.7,2.4)--node[lab] {$\CB(g,t^*)$} (9.9,2.85);
  \draw[msg,draw=rrRed] (5.7,2.4)--(14.1,2.85);
  \node[state,draw=rrRed,minimum width=8.0cm] at (10.0,3.55)
    {$Q$ and $L'$ also freeze generation $g$.\\
     Correct freeze times differ by at most $\Delta$.};
  \node[lab] at (7.8,4.35)
    {All preserve completed commits and process old votes/QCs; no new DATA votes.};

  \node[phase] at (.1,5.0) {(ii) Fix local locks and send STATUS to the predetermined leader};
  \node[wait] at (5.7,5.8) {Wait $\Delta$\\from $a_P$};
  \node[wait] at (9.9,5.8) {Wait $\Delta$\\from $a_Q$};
  \node[wait] at (14.1,5.8) {Wait $\Delta$\\from $a_{L'}$};
  \node[state,minimum width=12.6cm] at (7.9,6.75)
    {At own $a_r+\Delta$: fix highest-QC entry lock $Q_r$; enter $(g+1,0)$.\\
     Each STATUS binds the new view and carries $Q_r$ with complete verifying history.};
  \draw[msg] (5.7,7.6)--node[lab,pos=.5] {$\mathrm{STATUS}(g+1,0,Q_P)$} (14.1,8.0);
  \draw[msg] (9.9,8.2)--node[lab,pos=.5] {$\mathrm{STATUS}(g+1,0,Q_Q)$} (14.1,8.65);
  \node[wait,minimum width=8.7cm] at (10.0,9.35)
    {$L'$ waits $2\Delta$ from its own entry.\\
     This covers entry skew and delivery of correct STATUS reports.};

  \node[phase] at (.1,10.2) {(iii)--(iv) NEWVIEW, anchor re-certification, and the first fresh child};
  \draw[msg] (14.1,10.8)--node[lab,pos=.53]
    {$\mathrm{NEWVIEW}(g+1,0,Q^*)$\\highest valid received QC} (5.7,11.25);
  \node[state,minimum width=12.6cm] at (7.9,12.0)
    {Validate $Q^*$ and history; check $\operatorname{rank}(Q^*)\ge\operatorname{rank}(Q_r)$.\\
     Check NEWVIEW equivocation; relay NEWVIEW and vote for its unchanged anchor $A$.};
  \draw[msg] (5.7,12.8)--node[lab] {ANCHOR votes to all} (14.1,13.25);
  \node[state,draw=rrAmber,minimum width=12.6cm] at (7.9,13.9)
    {Form $\QC_{g+1,0}(A)$ from $f+1$ ANCHOR votes.\\
     Re-certification starts \textbf{no commit timer}.};
  \draw[msg] (14.1,14.7)--node[lab,pos=.53]
    {$\mathrm{PROPOSE}(Y,\QC_{g+1,0}(A))$} (5.7,15.15);
  \node[state,draw=rrTeal,minimum width=12.6cm] at (7.9,15.9)
    {First fresh child $Y$: DATA vote and start $2\Delta$ timer.\\
     Timed commitment of $Y$ also commits any uncommitted ancestors, including $A$.};
  \node[anchor=west,align=left,font=\fontsize{9}{10.5}\selectfont] at (.1,16.95)
    {Time runs downward; arrows are representative. L1's finalized $L(g+1,0)$ is already known.\\
     Late old-generation proofs do not trigger another recovery; future proposals wait for local entry.};
\end{tikzpicture}

\caption{Generation-wide exit followed by SH-style recovery. Different active terms share one frozen generation and one L1-scheduled target $(g+1,0)$. ANCHOR re-certification starts no timer.}
\label{fig:recovery}
\end{figure}

\begin{figure}[htbp]\centering
% Native TikZ state diagram. L1 assigns identity, not an anchor or SMR commit.
\begin{tikzpicture}[
  x=1cm,y=1cm,font=\fontsize{9.5}{11}\selectfont,
  flow/.style={-{Stealth[length=1.8mm]},semithick},
  state/.style={draw=rrBlue!65,fill=rrBlue!3,rounded corners=3pt,
    align=center,inner sep=5pt,text width=4.5cm,minimum height=1.15cm},
  lab/.style={fill=white,inner sep=3pt,align=center},
  note/.style={align=center,text width=14.8cm,inner sep=3pt}
]
  \path[use as bounding box] (0,-10.8) rectangle (15.8,.75);
  \node[state,draw=black!45,fill=black!3,text width=13.2cm] (lone) at (7.9,0)
    {Finalized L1 configuration supplies the immutable schedule $L(g,t)$.\\
     Schedule verification precedes entry; there is no recovery transaction to L1.};

  \node[state] (active) at (3.0,-2.25)
    {Active $(g,t)$\\Fresh proposals and DATA votes;\\$2\Delta$ commit timers};
  \node[state,draw=rrTeal,fill=rrTeal!4] (next) at (12.7,-2.25)
    {Active $(g,t+1)$\\Same generation and schedule;\\fresh height range advances};
  \draw[flow,draw=rrTeal] (active)--node[lab]
    {Terminal DONE certificate\\$\DC_{g,t}$ gates entry and votes} (next);
  \draw[flow,dashed,draw=black!45] (lone.south west)--(active.north);
  \draw[flow,dashed,draw=black!45] (lone.south east)--(next.north);

  \node[state,draw=rrRed,fill=rrRed!4] (frozen) at (3.0,-5.2)
    {Frozen generation $g$\\Cancel all its pending timers;\\retain completed commits};
  \draw[flow,draw=rrRed] (active)--node[lab,align=left,pos=.51]
    {Any complete valid\\current-generation $\CB(g,t^*)$} (frozen);
  \draw[flow,draw=rrRed] (next.south)--(12.7,-5.2)--node[lab,pos=.53]
    {A proof from an earlier\\term still freezes $g$} (frozen.east);

  \node[state,draw=rrAmber,fill=rrAmber!4] (recovery) at (3.0,-8.2)
    {Recover into $(g+1,0)$\\Wait $\Delta$; set lock;\\send STATUS to $L(g+1,0)$};
  \draw[flow,draw=rrAmber] (frozen)--node[lab,align=left,pos=.51]
    {One recovery per $g$;\\target independent of $t^*$} (recovery);
  \node[state] (newactive) at (12.7,-8.2)
    {Active $(g+1,0)$\\NEWVIEW anchor\\re-certified;\\fresh children start timers};
  \draw[flow] (recovery)--node[lab]
    {$L(g+1,0)$ waits $2\Delta$;\\NEWVIEW;\\ANCHOR re-certification} (newactive);
  \node[note] at (7.9,-10.1)
    {L1 selects the leader identity; the leader proposes an anchor subject to each replica's lock.\\
     Previous-generation CBs cannot advance the new generation.};
\end{tikzpicture}

\caption{The state transitions and L1 boundary. L1 supplies leader identities in advance; local evidence triggers recovery and SH entry locks constrain the proposed anchor.}
\label{fig:states}
\end{figure}

\FloatBarrier


\section{Safety}
\label{sec:safety}
We first establish the timing and history invariants, then prove preservation
across generations. All statements refer to admissible certificates; a
certificate's signatures alone are insufficient if its entry history is absent.

\begin{lemma}[Entry skew and recovery target]
\label{lem:entry}
For each entered generation, correct initial-view entry times differ by at
most $\Delta$. If its first correct freeze is at $a$, every correct replica
freezes that same generation by $a+\Delta$ and enters the fixed successor
generation at a time in $[a+\Delta,a+2\Delta]$. Different triggering terms
or witnesses cannot change the target.
\end{lemma}
\begin{proof}
The initial claim is \cref{ass:input}. For the induction step, immediately
relaying the complete generation-$g$ proof delivers it within $\Delta$.
No correct replica can have completed generation-$g$ recovery before the
first correct freeze $a$. A replica still finishing generation $g-1$
buffers the proof: by the inductive entry-skew bound it enters $g$ no
later than $\Delta$ after the first correct generation-$g$ entry, hence no
later than $a+\Delta$. It processes the buffered proof on entry. Thus
$a_r\in[a,a+\Delta]$. The local wait gives
$e_r=a_r+\Delta\in[a+\Delta,a+2\Delta]$. All proofs name generation $g$,
and the transition function always returns $(g+1,0)$ with the same
pre-finalized leader. A buffered next-generation proof may immediately
freeze that generation after entry; it does not change the entry time or
the completed predecessor transition.
\end{proof}

\begin{lemma}[Canonical history or signed evidence]
\label{lem:fork}
Two admissible histories in one generation have the same opening and are
prefix-comparable, or their retained signed material supplies a valid quit
proof for that generation.
\end{lemma}
\begin{proof}
Different opening blocks supply the two valid NEWVIEW signatures of
\cref{sec:evidence}(c). Generation zero has its fixed genesis opening.
For the same opening, fresh height determines the term and scheduled
leader. At the first divergent height, both histories therefore contain
different DATA proposals signed by the same leader at the same
$(g,t,k)$. Both entry paths are complete, so these proposals give case (b).
Terminal commit certificates are required in both paths at each normal
boundary. No unverifiable future-term proposal is used in this argument.
\end{proof}

\begin{lemma}[Protection of a completed DATA timer]
\label{lem:timer}
Suppose a correct replica DATA-votes for $B$ at time $s$ in generation $g$
and its own timer commits $B$ at $s+2\Delta$. Every correct replica knows
its DATA QC by that time. No admissible generation-$g$ QC conflicting with
$B$ can exist, including a QC assembled later from old votes. No ANCHOR QC
under a different generation-$g$ opening can exist.
\end{lemma}
\begin{proof}
No correct replica freezes generation $g$ by $s+\Delta$: its immediately
relayed proof would reach the committer by $s+2\Delta$ and cancel the timer.
The complete proposal forwarded at $s$ reaches every correct replica by
$s+\Delta$. A replica still finishing prior recovery has entered $g$ by
then, by \cref{lem:entry}; a replica in a later generation would require
an earlier freeze, already excluded.

An incompatible retained admitted history would expose evidence by
\cref{lem:fork}, cause a freeze by $s+\Delta$, and contradict successful
commitment. Otherwise the recipient catches up through the supplied
opening and terminal commit certificates. If it has passed $B$ on this
certified history, its ancestry contains $QC(B)$ and it relays that QC.
One such correct relay supplies the QC everywhere by $s+2\Delta$.
If none has passed $B$, all correct replicas have voted or now vote for
$B$ by $s+\Delta$. There are at least $f+1$ correct votes; their multicast
supplies the QC everywhere by expiry. Votes remain processable even after
a subsequent freeze.

For a conflicting QC, choose a correct signer. If it processes the
committer's proposal before attempting the conflicting vote, the retained
fork makes it freeze instead of vote. Otherwise that vote was sent by
$s+\Delta$, together with its conflicting signed proposal or NEWVIEW.
The relay reaches the committer by $s+2\Delta$, cancelling the timer.
This is a contradiction in either case. The same argument applies to an
ANCHOR vote for any different opening, even when the opening blocks are
comparable: the two signed openings alone suffice to freeze. It concerns
the time of a correct vote, not the later time at which Byzantine nodes
assemble or disclose a QC.
\end{proof}

\begin{lemma}[Locks inherit completed timers]
\label{lem:locks}
If a correct timer commits $B$ in generation $g$, every correct replica's
entry lock for every later generation ranks at least as high as its
protecting DATA QC.
\end{lemma}
\begin{proof}
Let the timer start at $s$ and let $a$ be the first correct freeze of $g$,
if there is one. \Cref{lem:timer}'s cancellation argument gives
$a>s+\Delta$. Every entry-lock snapshot occurs at
$a_r+\Delta\ge a+\Delta>s+2\Delta$, after every correct replica knows the
protecting QC. Its highest-QC lock therefore ranks at least that high.
Subsequent recovery preserves nondecreasing lock ranks. If $g$ never
freezes, no later generation is entered and the assertion is vacuous.
\end{proof}

\begin{lemma}[Unique extensibility]
\label{lem:extend}
If a correct timer commits $B$, every admissible QC of rank at least
$\rank(QC(B))$ certifies a block extending $B$.
\end{lemma}
\begin{proof}
Within $B$'s generation, \cref{lem:timer} excludes alternative certified
openings and incompatible QCs. Under the one certified opening, canonical
term ranges increase with both term and fresh height. A QC of at least
$B$'s rank therefore has height at least $h(B)$; it cannot be a strict
ancestor, and hence must extend $B$.

Suppose a later-generation QC of sufficient rank fails to extend $B$.
Choose one with the least rank among such QCs. If it is an ANCHOR QC for
$A$, a correct ANCHOR voter accepted an older-generation justification
$Q_A$ whose rank is at least its entry lock, hence at least $QC(B)$ by
\cref{lem:locks}. Since $A$ does not extend $B$, $Q_A$ belongs to the
same set and has strictly smaller rank, a contradiction.
If the chosen QC is a DATA QC, its parent QC is in the same later
generation, including the ANCHOR parent of its first fresh child and the
preceding-term terminal parent at ordinary handoff. Its rank is strictly
smaller but still greater than $QC(B)$. A parent extending $B$ would make
its child extend $B$, so this parent also belongs to the set, again a
contradiction. The lexicographic ranks are well-ordered, and equal-rank
distinct blocks do not affect either strict-decrease argument.
\end{proof}

\begin{lemma}[Terminal commit certificates]
\label{lem:done}
Every valid $\DC_{g,t}(T)$ witnesses a correct completed DATA timer for
$T$. Importing its committed prefix and using it for normal admission
preserves safety. In particular, an admissible quit proof for term $t>0$
already witnesses a fresh commitment in that generation.
\end{lemma}
\begin{proof}
Among $f+1$ distinct DONE signers at least one is correct. That signer
issued DONE only after its own terminal DATA timer succeeded; imported
certificates cannot cause it to sign DONE. The certificate therefore
inherits a genuine timer commitment. A later term's admission contains
its predecessor terminal commit certificate, whose block is fresh above
the generation opening. This applies even if the proof arrives before
the recipient has locally entered that term.
\end{proof}

\begin{theorem}[SMR safety and validity]
\label{thm:safety}
All correct committed logs are prefix-comparable, completed commitments
are never withdrawn, and every committed block satisfies application
validity. All correct exits from generation $g$ target $(g+1,0)$ and its
unique L1-scheduled leader.
\end{theorem}
\begin{proof}
For two direct timer commitments, order their protecting QCs by rank.
\Cref{lem:extend} says the higher-ranked certified block extends the
lower-ranked committed block. Indirect ancestor commitments inherit this
comparability. Every imported terminal commitment inherits a direct timer
commitment by \cref{lem:done}, so imports introduce no additional
incompatible log. All accepted proposals and certificate histories satisfy
application validity; their committed ancestors do as well.
The protocol cancels unfinished timers on a freeze and never revokes
completed commits. \Cref{lem:entry} and the unique finalized schedule
give the recovery-target and leader assertions.
\end{proof}

\section{Liveness}
\label{sec:liveness}
\begin{lemma}[A correct leader's progress]
\label{lem:honest}
An unfrozen correctly entered term with a correct leader, and no newly
delivered earlier-term quit proof, gives every correct replica at least
$p$ distinct admitted fresh positions, evidenced by its own DATA vote or
a valid DATA QC, by $(2p+4)\Delta$ after that replica's entry,
for $1\le p\le K$. A correct initial leader cannot be quit before its
generation has made a fresh commitment.
\end{lemma}
\begin{proof}
For a recovered initial term let $e$ be the earliest correct entry.
The leader enters by $e+\Delta$ and finishes its own $2\Delta$ wait by
$e+3\Delta$. All correct STATUS reports have arrived by then. Its chosen
QC rank is at least every fixed entry lock, so its one valid NEWVIEW
reaches all correct replicas by $e+4\Delta$ and triggers their ANCHOR
votes. The leader has an ANCHOR QC by $e+5\Delta$, and its first fresh
proposal is processed by $e+6\Delta$.
Each subsequent proposal/QC iteration costs at most $2\Delta$; the
$p$th proposal is processed by $e+(2p+4)\Delta$. A recipient either
votes or already has QC-backed catch-up evidence for that position.
Thus its progress count is at least $p$. Every local entry is at least
$e$, giving the stated local deadline. Generation zero omits the anchor
exchange and is faster.

Normal entry times differ by at most $\Delta$: the first correct receipt
of a valid terminal commit certificate relays it immediately, unless the
generation freezes. Its new correct leader proposes the first child
immediately on entry (or has pre-sent it). First votes arrive within
$2\Delta$ of the earliest normal entry, and later votes retain the
$2\Delta$ iteration bound, which is within the stated deadlines.
At exact deadlines messages are processed before BLAME callbacks.

A correct leader issues one opening and one proposal per slot, so it
does not supply equivocation evidence. The timely correct votes prevent
correct BLAME signatures; Byzantine signatures alone cannot reach $f+1$.
Before a first generation commitment there is also no admissible
higher-term quit proof, because it requires a terminal commit certificate
by \cref{lem:done}. Thus a correct initializer cannot be frozen before
fresh progress. Its terminal timer succeeds, or a higher-term proof
already witnesses that success.
\end{proof}

\begin{lemma}[No-progress generations terminate or commit]
\label{lem:no-progress}
If a generation makes no fresh correct timer commitment, it cannot have
an admissible normal entry or higher-term quit proof. If its initial
leader is correct it makes fresh progress; otherwise a permanent stall
eventually produces a complete initial-term BLAME certificate.
\end{lemma}
\begin{proof}
Normal entry and higher-term proof admission require DONE, which is
impossible without a fresh timer commitment. Only term zero can be active
without progress. The correct-initializer conclusion is
\cref{lem:honest}. If fresh DATA votes stop before quota completion, each
correct replica eventually reaches a deficient progress test for this
same term, issues its one complaint, and receives the complaints of at
least $f+1$ correct replicas. They assemble and relay a complete quit
proof. If instead all fresh DATA votes occur, their QCs arrive, and their
timers commit unless a quit proof freezes the generation; that proof
already starts finite recovery. Equivocation likewise immediately starts
recovery once exposed. Infinite silence cannot leave a no-progress
generation active forever.
\end{proof}

\begin{theorem}[Committed-height progress]
\label{thm:height}
Under \cref{ass:network,ass:l1,ass:input} and initializer fairness,
every correct replica eventually commits arbitrarily large heights.
\end{theorem}
\begin{proof}
Suppose first that infinitely many generations are entered. A valid
recovery transition finishes after its local wait; no additional agreement
or L1 transaction can stall entry. By initializer fairness, infinitely
many of these generations have correct initial leaders. Each such
generation completes its initial terminal commitment by
\cref{lem:honest}. Its opening extends every older correct commitment:
any correct ANCHOR vote accepts an older QC at least as high as its
entry lock, and \cref{lem:extend,lem:locks} preserve all earlier timer
commits. Its fresh terminal height is therefore strictly above the
previous committed maximum. The terminal's protecting QC reaches every
correct replica, and subsequent recovery locks preserve it. Thus direct
commit heights grow without bound.

Now suppose there are finitely many entered generations. Their last
generation cannot freeze, since a freeze would cause another finite
transition. If normal terms advance infinitely, their admission
certificates witness terminal commitments at increasing canonical heights.
If only finitely many terms advance, all correct replicas eventually
catch up through the disclosed entry certificates. In their last common
term, either insufficient DATA progress produces $f+1$ correct BLAME
signatures and recovery, or all $K$ blocks are voted and certified.
With no freeze their timers commit, all correct DONE signatures are
multicast, and the terminal commit certificate advances the term. Both
possibilities contradict a permanently last term.

Finally take a direct commitment at any height. Its protecting QC is
known everywhere and all later accepted openings extend it. Every
correct replica either imports a terminal commitment or commits a later
DATA descendant on its own timer. In an infinite-generation execution,
each correct-initializer term reaches its terminal position. Every correct
replica either casts that terminal DATA vote or has already imported its
terminal DONE. The resulting timers finish unless a later-term quit proof
carries that terminal DONE, which is relayed and imported by all. Hence
arbitrarily high committed prefixes are eventually
installed at every correct replica, not merely at one witness.
\end{proof}

\begin{theorem}[Request liveness]
\label{thm:requests}
With both schedule fairness properties and the input service of
\cref{ass:input}, every persistently admissible correct-client request is
eventually committed by all correct replicas.
\end{theorem}
\begin{proof}
Infinitely many recoveries give infinitely many committed initial terms
with correct leaders, as in \cref{thm:height}. Finitely many recoveries give
an infinite unfrozen generation with infinitely many completed correct
terms by normal-term fairness. Thus there are infinitely many completed
correct terms in either case.

Fix a request after its dissemination to all correct replicas. It has
finitely many predecessors in the common admission order. At every
completed correct term, its leader's first fresh block contains the
oldest remaining requests that fit, or those requests already occur in
the ancestry committed by that term. Speculative inclusion never removes
a request from the pending queue. Induction over the finite predecessors
therefore places the fixed request in a completed correct term's prefix.
Safety and committed-prefix dissemination then deliver it to every
correct replica.
\end{proof}

\section{Complexity and performance}
\label{sec:complexity}
The counts below use ordinary signatures, explicit certificates, and
deduplication as in \cref{sec:processing}. Let $\sigma$ be signature size,
$b$ the fresh block size, and $A$ the other ancestry/admission bytes
actually carried with a representative proposal or recovery report.
These are transmitted sizes; a cached dependency has zero retransmission
cost only where the receiver can validate its reference.

\begin{center}\small
\begin{tabularx}{\linewidth}{@{}p{.20\linewidth}p{.19\linewidth}Y@{}}
\toprule
Operation & Correct messages & Conservative bytes with full forwarding\\
\midrule
Fresh DATA block & $O(n^2)$ &
$O(n^2(b+A+n\sigma))$: proposals and their parent QC are forwarded;
votes themselves cost $O(n^2\sigma)$.\\
Terminal DONE & $O(n^2)$ per term &
$O(n^2\sigma)$ for shares; forwarding explicit DONE certificates can
add $O(n^3\sigma)$, with entry-history bytes charged separately.\\
Failure recovery control & $O(n^2)$ per generation &
$O(n^2(A+n\sigma))$: STATUS to one leader, NEWVIEW forwarding and
all-to-all ANCHOR votes; quit and commitment evidence are charged.\\
L1 schedule & Outside local execution &
One finalized configuration/schedule witness, amortized separately.
No per-block or per-recovery transaction.\\
\bottomrule
\end{tabularx}
\end{center}
The $O(n^2)$ message order agrees with raw standard-synchronous SH.
Late DATA votes, DATA QCs, and DONE traffic processed while frozen are
charged to their originating blocks and terms, not counted again as
bounded new recovery control traffic.
The new DONE traffic contributes $O(n^2/K)$ messages per fresh block
amortized over a $K$-block term. Without certificate compression, message
complexity and bit complexity are different. Unbounded distinct Byzantine
traffic is outside this accounting.

\paragraph{Block and recovery latency.}
An admitted DATA vote starts the same $2\Delta$ timer as SH.
An ordinary unbuffered proposal under homogeneous actual one-hop delay
$\delta\le\Delta$ commits after $\delta+2\Delta$ at a follower if its
timer survives. For a speculative proposal, physical send-to-commit also
includes the admission-buffering interval; the $2\Delta$ bound starts at
the DATA vote, not at its early transmission.
From earliest freeze $a$, all enter the successor generation by
$a+2\Delta$. A correct new leader sends NEWVIEW by $a+4\Delta$, obtains
the ANCHOR QC by $a+6\Delta$, and reaches all correct replicas with its
first fresh child by $a+7\Delta$. All that child's timers expire by
$a+9\Delta$ if not frozen. With commit-gated admission, any admissible
higher-term interruption already witnesses a terminal commitment.
The two waiting stages are $\Delta$ and $2\Delta$, as in SH; they are
not the whole end-to-end recovery time.

\paragraph{Term boundary and throughput.}
Within one term, QCs pipeline proposals at network speed, approximately
one block per $2\delta$ in the homogeneous good case. The terminal commit
gate puts a conservative wait on the ordinary handoff path. Let $c_t$
be the maximum, over correct replicas, of their first local processing
times for term $t$'s first admitted fresh position, by a DATA vote or
QC-backed catch-up.
If consecutive terms have correct leaders and no freeze, then
\[
 c_{t+1}-c_t\ \le\ 2\Delta+(2K-1)\delta
\]
with a buffered speculative successor, and the conservative bound without
speculation is $2\Delta+2K\delta$. If a correct replica already knows
terminal DONE, its prompt relay supplies it everywhere within $\delta$.
Otherwise the first-position QC reaches the leader by $c_t+\delta$:
either all correct first votes have been sent, or a QC-backed recipient
relays the QC. For $K\ge2$, later proposal/QC iterations reach the terminal
position by $c_t+(2K-2)\delta$. Each correct recipient then votes for that
terminal block, or has already passed it using DONE and relays the
certificate. In the former case terminal timers finish within another
$2\Delta$ and DONE shares arrive within one more $\delta$; the latter
case is faster. For $K=1$, all first-position votes have occurred by
$c_t$, or completed-term catch-up already supplies DONE.
The speculative first proposal is already available by this DONE bound.
Ignoring bandwidth and input queues, an uninterrupted sequence of correct
terms consequently supports at least
\[
 \frac{K}{2\Delta+(2K-1)\delta}
\]
fresh blocks per time unit under this timing envelope.
This is a conservative good-case service bound, not an exact rate for
arbitrary Byzantine participation or link delays. The handoff cost
amortizes as $K$ increases; it is not identical to SH's stable-leader
throughput.

\section{Immediate-QC rotation as a further optimization}
\label{sec:optimization}
The preferred faster variant would admit $(g,t+1)$ on $N_{g,t}$ alone.
Its proposal and DATA-vote pipeline can cross a leader boundary before
the predecessor's timers expire. This manuscript proves the DONE-gated
variant; it does not silently assign that faster rule the same liveness theorem.

To see the missing obligation, take $n=3$, $f=1$, $K=1$, with a faulty
leader after short correct terms and actual delay $\delta\ll\Delta$.
Correct terms can form QCs and hand off before any $2\Delta$ timer expires.
The faulty successor equivocates, freezes the generation, and cancels
all those timers. SH's recovery re-certifies an inherited anchor without
a timer. A fair L1 schedule can repeat this pattern in subsequent
generations: certified height grows while committed height stops.
Agreement on the next leader does not resolve this execution.

Immediate-QC rotation therefore needs another mechanism that converts
certified progress into irreversible commitment despite repeated recovery,
with a separate preservation and progress proof. A recovery anchor cannot
be committed merely because it has $f+1$ votes. The current commit gate
is the explicit conservative choice, and the faster variant remains a
separate design objective.

\section{Changes relative to original SH}
\label{sec:diff}
\begin{center}\small
\begin{tabularx}{\linewidth}{@{}p{.22\linewidth}Y Y@{}}
\toprule
Step & Original standard-synchronous SH & RR v2\\
\midrule
View and leader & Scalar view and cyclic leader. &
$v=(g,t)$; pre-finalized $L(g,t)$.\\
Normal production & Stable leader until failure. &
$K$ fresh positions per term; terminal DONE gates successor DATA votes.\\
Proposal pipeline & Parent QC enables a child. &
Same within a term; early successor proposals are buffered until admission.\\
Quit & Cancel current-view timers. &
One complete current-generation proof cancels all generation timers.\\
STATUS & After $\Delta$, lock and report highest QC to next leader. &
Same operation; every trigger targets $(g+1,0)$ independently of old $t$.\\
NEWVIEW & Leader waits $2\Delta$ and chooses highest QC. &
Same waits and rank rule; one immutable opening fixes quota origins.\\
First vote & Re-certify inherited anchor without timer. &
Same ANCHOR re-vote; no BA-selected anchor or recovery commit.\\
\bottomrule
\end{tabularx}
\end{center}
The recovery communication pattern and original DATA timer are retained.
The additional public normal-entry certificate is essential to this
version's liveness proof and is visibly marked in the protocol.

\begin{thebibliography}{2}\small
\bibitem{sync}
I. Abraham, D. Malkhi, K. Nayak, L. Ren, and M. Yin.
\emph{Sync HotStuff: Simple and Practical Synchronous State Machine
Replication}. IEEE Symposium on Security and Privacy, 2020, Section II,
Figures 1--2.
\href{https://conferences.computer.org/sp/pdfs/sp/2020/349700a106.pdf}{Published paper}.
\bibitem{schneider}
F. B. Schneider. Implementing fault-tolerant services using the state
machine approach: A tutorial. \emph{ACM Computing Surveys} 22(4),
299--319, 1990. \href{https://doi.org/10.1145/98163.98167}{DOI}.
\end{thebibliography}
\end{document}
